%! Author = sriadityavedantam
%! Date = 4/7/26

\section{Singular and Non-singular Points}

Let $C : \{f(x,y) = 0\}\subseteq \a^2$, and let $P\in C$ such that $f(P) = 0$.

\begin{definition}[Singular Point]
    A point $P$ is singular if and only if
    \[
        \frac{\partial f}{\partial x}(P) = 0
        \quad\text{and}\quad
        \frac{\partial f}{\partial y}(P) = 0.
    \]
\end{definition}

\begin{lemma}{}{
    If $P = (0,0)\in C$, then $f(x,y)$ has no constant term.
}
\end{lemma}

\begin{lemma}{}{
    If $P$ is singular, then $f(x,y)$ has no linear term.
}
\end{lemma}

\begin{definition}[Smooth Curve]
    If for all $P\in C$, the point $P$ is non-singular, then $C$ is smooth.
\end{definition}

Smoothness is invariant under affine change of coordinates.
It suffices to check that
\[
    x^2 - y^2 - 1 = 0
    \quad\text{and}\quad
    x^2 - y = 0
\]
are smooth.

\begin{theorem}{}{
    An irreducible curve $C$ has only finitely many singular points.
}
Let $P\in C$.
Define
\[
    C' \coloneqq \left\{ \frac{\partial f}{\partial x} = 0 \right\}.
\]
If $P$ is singular, then $P\in C\cap C'$.
From the lemma, either $C\cap C'$ is finite or $C$ divides $C'$, that is, $f \mid \frac{\partial f}{\partial x}$.
However,
\[
    \deg\!\left(\frac{\partial f}{\partial x}\right) < \deg(f),
\]
so this is impossible unless $\frac{\partial f}{\partial x} = 0$.
Similarly, $\frac{\partial f}{\partial y} = 0$.
If both partial derivatives are zero, then $f$ is constant.
Since we are working over $\c$ with $\char = 0$, this cannot define a curve.
Hence there are only finitely many singular points.
\end{theorem}

\begin{definition}[Multiplicity]
    Let $C$ be a curve and $P\in C$.
    Assume $P=(0,0)$.
    Then $\Mult(P)$ is the smallest degree of a nonzero monomial appearing in $f(x,y)$.
\end{definition}

\begin{align*}
    f(x,y) &= x - y + x^2,\\
    &\implies \Mult(0,0) = 1.
\end{align*}

Recall that $\Mult(P) = 1$ if and only if $f$ contains a linear term, which implies that $C$ is smooth at $P$.

\begin{align*}
    f(x,y) &= xy,\\
    &\implies \Mult(0,0) = 2.
\end{align*}

\begin{align*}
    f(x,y) &= y^2 - x^3,\\
    &\implies \Mult(0,0) = 2.
\end{align*}

If $C$ is a curve of degree $n$, then any point on $C$ has multiplicity at most $n$.